\documentclass[11pt,a4paper]{article}

\usepackage{amsmath, amssymb, amsthm}
\usepackage[margin=1in]{geometry}
\usepackage{hyperref}

\hypersetup{
    colorlinks=true,
    linkcolor=blue,
    urlcolor=blue
}

\newtheorem*{exercise}{Exercise}
\newtheorem*{solution}{Solution}

\title{\textbf{LADR 4e --- Alternative Solutions}}
\author{vecljox}
\date{}

\begin{document}

\maketitle

\noindent
This document collects alternative proofs for selected exercises in
\textit{Linear Algebra Done Right} (4th edition) by Sheldon Axler.
Each solution here differs from the official or standard approach
and aims to be more direct or illuminating.

\bigskip
\noindent
Licensed under CC BY-NC-SA 4.0.

\bigskip
\hrule
\vspace{1em}

\section*{Chapter 6.B}

\begin{exercise}[6.B.9]
Suppose $e_1, \dots, e_m$ is the result of applying the Gram--Schmidt
procedure to a linearly independent list $v_1, \dots, v_m \in V$.
Prove that $\langle v_k, e_k \rangle > 0$ for each $k = 1, \dots, m$.
\end{exercise}

\begin{solution}
By the Gram--Schmidt procedure, $e_k = \dfrac{f_k}{\|f_k\|}$, so
\[
    \langle v_k, e_k \rangle
    = \frac{\langle v_k, f_k \rangle}{\|f_k\|}.
\]
It suffices to show that $\langle v_k, f_k \rangle > 0$.

Substituting the definition
\[
    f_k = v_k - \sum_{i=1}^{k-1}
    \frac{\langle v_k, f_i \rangle}{\|f_i\|^2}\, f_i
\]
and using $f_i = \|f_i\|\, e_i$, we obtain
\[
    \langle v_k, f_k \rangle
    = \|v_k\|^2 - \sum_{i=1}^{k-1} |\langle v_k, e_i \rangle|^2.
\]
By Bessel's inequality,
\[
    \sum_{i=1}^{k-1} |\langle v_k, e_i \rangle|^2
    \leq \|v_k\|^2,
\]
so $\langle v_k, f_k \rangle \geq 0$.

If equality holds, then
$\|v_k\|^2 = \sum_{i=1}^{k-1} |\langle v_k, e_i \rangle|^2$.
By Exercise~6.B.3, this implies
$v_k \in \operatorname{span}(e_1, \dots, e_{k-1})
= \operatorname{span}(v_1, \dots, v_{k-1})$,
contradicting the linear independence of $v_1, \dots, v_m$.

Therefore $\langle v_k, f_k \rangle > 0$, and hence
\[
    \langle v_k, e_k \rangle
    = \frac{\langle v_k, f_k \rangle}{\|f_k\|} > 0.
    \qedhere
\]
\end{solution}

\section*{Chapter 7.F}

\begin{exercise}[7.F.3]
Suppose $T \in \mathcal{L}(V, W)$ and $v \in V$. Prove that
\[
    \|Tv\| = \|T\|\,\|v\|
    \quad \Longleftrightarrow \quad
    T^*Tv = \|T\|^2 v.
\]
\end{exercise}

\begin{solution}
($\Leftarrow$) Suppose $T^*Tv = \|T\|^2 v$. Taking the inner product
of both sides with $v$ gives
\[
    \langle T^*Tv, v \rangle = \langle \|T\|^2 v, v \rangle.
\]
The left-hand side equals $\langle Tv, Tv \rangle = \|Tv\|^2$, and the
right-hand side equals $\|T\|^2 \|v\|^2$. Hence
$\|Tv\|^2 = \|T\|^2 \|v\|^2$, and taking square roots yields
$\|Tv\| = \|T\|\,\|v\|$.

\medskip
($\Rightarrow$) Suppose $\|Tv\| = \|T\|\,\|v\|$. Let
$S = \|T\|^2 I - T^*T$. We first show that $S$ is a positive operator.
For each $u \in V$,
\[
    \langle Su, u \rangle
    = \|T\|^2 \|u\|^2 - \langle T^*Tu, u \rangle
    = \|T\|^2 \|u\|^2 - \|Tu\|^2
    \geq 0,
\]
where the inequality follows from $\|Tu\| \leq \|T\|\,\|u\|$. Since
$T^*T$ and $I$ are self-adjoint, so is $S$; thus $S$ is positive.
Next, using the hypothesis $\|Tv\| = \|T\|\,\|v\|$,
\[
    \langle Sv, v \rangle
    = \|T\|^2 \|v\|^2 - \|Tv\|^2
    = 0.
\]
By 7.43, a positive operator $S$ with $\langle Sv, v \rangle = 0$
satisfies $Sv = 0$. Therefore $\|T\|^2 v - T^*Tv = 0$, that is,
\[
    T^*Tv = \|T\|^2 v. \qedhere
\]
\end{solution}

\begin{exercise}[7.F.8]
\begin{enumerate}
    \item[(a)] Prove that if $T \in \mathcal{L}(V)$ and $\|I - T\| < 1$,
    then $T$ is invertible.
    \item[(b)] Suppose that $S \in \mathcal{L}(V)$ is invertible. Prove
    that if $T \in \mathcal{L}(V)$ and
    \[
        \|S - T\| < \frac{1}{\|S^{-1}\|},
    \]
    then $T$ is invertible.
\end{enumerate}
\end{exercise}

\begin{solution}
We prove (b) directly; part (a) follows as the special case $S = I$.

Since $V$ is finite-dimensional, it suffices to show that $T$ is
injective, i.e., $Tv \neq 0$ whenever $v \neq 0$.

We first observe that $S$ admits a uniform lower bound on its stretching:
for every $v \in V$,
\[
    \|v\| = \|S^{-1} S v\| \leq \|S^{-1}\|\, \|Sv\|,
\]
which rearranges to
\[
    \|Sv\| \geq \frac{1}{\|S^{-1}\|}\, \|v\|. \tag{$\ast$}
\]

Now let $v \in V$ with $v \neq 0$. Writing $T = S - (S - T)$ gives
\[
    Tv = Sv - (S - T)v.
\]
By the reverse triangle inequality,
\[
    \|Tv\| \geq \|Sv\| - \|(S - T)v\|.
\]
Applying ($\ast$) to the first term and $\|(S - T)v\| \leq \|S - T\|\,\|v\|$
to the second,
\[
    \|Tv\|
    \geq \frac{1}{\|S^{-1}\|}\,\|v\| - \|S - T\|\,\|v\|
    = \left( \frac{1}{\|S^{-1}\|} - \|S - T\| \right) \|v\|.
\]
By hypothesis $\|S - T\| < \dfrac{1}{\|S^{-1}\|}$, so the coefficient
on the right is strictly positive. Since $\|v\| > 0$, we conclude
$\|Tv\| > 0$, hence $Tv \neq 0$.

Therefore $T$ is injective, and thus invertible. \qedhere
\end{solution}

\end{document}
